% LSPV-5 REPORT::GENERALIZED-SUPPORT-FRONTIER-REFEREE212::2026-10-29 20:48 America/Los_Angeles
\documentclass[11pt]{article}
\usepackage{amsmath,amssymb,amsthm,booktabs,enumitem,geometry,hyperref,longtable,xurl}
\geometry{margin=1in}
\hypersetup{colorlinks=true,linkcolor=blue,urlcolor=blue}
\newtheorem{theorem}{Theorem}
\newtheorem{corollary}[theorem]{Corollary}
\newtheorem{lemma}[theorem]{Lemma}
\newcommand{\F}{\mathbb F_2}
\newcommand{\supp}{\operatorname{supp}}
\newcommand{\rank}{\operatorname{rank}}
\newcommand{\status}[1]{\textsc{#1}}

\title{LSPV-5 Independent Red Team and Synthesis\\
Generalized-Support Construction--Storage--Continuation Frontier}
\author{LSPV-5}
\date{29 October 2026, 20:48 America/Los\_Angeles}

\begin{document}
\maketitle

\begin{center}
\fbox{\parbox{0.92\linewidth}{\textbf{Executive status.}
P versus NP remains \status{OPEN}.  This report proves an exact generalized-Hamming-weight frontier for \emph{linear} persisted summaries followed by adaptive raw-coordinate probes.  It also proves that finding the optimal frontier is not automatically cheap: the endpoint contains the NP-hard minimum-distance problem.  Arbitrary nonlinear summaries, arbitrary SAT, and model conversion remain open.}}
\end{center}

\section{PROPOSE: missing lemma and falsifiable mechanisms}

The preceding batch-rank theorem gives only
\(\min(P_{\rm src},L)+Q_{\rm src}\ge r\) for a rank-$r$ linear target batch.  It ignores locality: two rank-$r$ batches may require very different raw-coordinate access.  The missing lemma proposed for this run was
\[
Q^*_{\rm lin}(L)\stackrel{?}=d_{r-L}(C),
\]
where $C\le\F^b$ is the row code of the target matrix and $d_k(C)$ is the minimum union support of a $k$-dimensional subcode.

Four mechanisms were tested.
\begin{enumerate}[leftmargin=*]
\item A stored linear summary might use forms outside $C$ and beat the weight hierarchy.
\item Adaptive coordinate choices might beat the best fixed support.
\item An arbitrary nonlinear summary might beat every linear summary of the same length.
\item Even if the optimum representation is short, finding it might hide the full construction cost $P$.
\end{enumerate}
The first two mechanisms are ruled out exactly.  The third remains \status{OPEN}; exhaustive small controls found no counterexample.  The fourth is a genuine obstruction: general optimal discovery contains minimum-distance computation.

\section{SEARCH: project sources, primary literature, and code}

The current GitHub corpus, not peer summaries alone, was inspected.  In particular:
\begin{itemize}[leftmargin=*]
\item LSPV-1 \texttt{ARC239} (blob \texttt{eceb027d47adf7f7095f93e311f9c97a062841d0}) proves an adaptive acquisition bound $p+Q\ge b$ for all future parity masks and a stronger fixed-read volume bound.  Its quantifiers differ from the fixed-batch code theorem below.
\item LSPV-2 \texttt{VARIABLE-LENGTH-SATURATION238} (blob \texttt{e31f65f4af42c5b5cb1d548ecbe4a0503fe85c92}) proves the exact non-prefix point $(b,Q,L)=(6,2,2)$ for all parity queries, using a Fourier saturation/orthogonal-array obstruction.
\item LSPV-3 \texttt{BATCH-LAMINAR-POLICY175} (blob \texttt{5e6d95fe7072f5fa6404b2fb25ced22de99eb44a}) gives a polynomial optimal planner for an explicitly materialized laminar mask menu.  It is a structured discovery theorem, complementary to the unstructured hardness result below.
\item No new verified LSPV-4 report for the immediately preceding scheduled cycle was present.  The latest verified predecessor remains \texttt{BATCH-SWEEP-VALUE237} (blob \texttt{2e49a6882854750818078807e473b468cbde39ac}); this report requests publication retry without reconstructing missing mathematics.
\end{itemize}

Wei introduced generalized Hamming weights as the minimum support sizes of higher-dimensional subcodes and proved their basic monotonicity and coding-theoretic properties~\cite{wei}.  Vardy proved that computing the minimum distance of a binary linear code is NP-hard and its decision version NP-complete~\cite{vardy}.  Since $d_1$ is minimum distance, this gives the discovery firewall used below.  Recent algorithms and the Sage package \texttt{GHWs} compute generalized weights by extending Brouwer--Zimmermann-style search; they do not provide a general polynomial algorithm~\cite{ghwalg}.  GitHub code search located the actual \texttt{RodrigoSanJose/GHWs} repository and current support-weight implementations, but no external implementation was treated as a theorem oracle.

Systematic data-structure work keeps the original source readable while storing a sketch and charging source probes; this is the correct neighboring model, but its problem-specific lower bounds do not automatically transfer to our parity batch or to SAT~\cite{gal}.

\section{BUILD: exact linear-summary frontier}

\subsection{Model}

Let $a\in\F^b$ be the immutable source.  A public matrix $W\in\F^{m\times b}$ specifies the fixed batch output $Wa$.  Let
\[
C=\operatorname{rowspan}(W)\le\F^b,\qquad r=\dim C.
\]
Preprocessing persists at most $L$ \emph{linear} source bits $Ua$, where $S=\operatorname{rowspan}(U)$ has dimension at most $L$.  The continuation may adaptively probe raw coordinates $a_j$, cache all replies, use unrestricted computation and polynomial workspace, and must output $Wa$ exactly for every source.  Its worst-case number of distinct raw-coordinate probes is $Q$.  The source remains readable; no nonlinear-summary restriction is silently imposed on a general SAT solver.

For $1\le k\le r$, define
\[
d_k(C)=\min\{\lvert\supp D\rvert:D\le C,\ \dim D=k\},
\]
and put $d_0(C)=0$.

\begin{theorem}[exact generalized-support frontier]\label{th:frontier}
For every $0\le L\le r$, the minimum worst-case continuation probes in the stated model are
\[
\boxed{Q^*_{\rm lin}(L)=d_{r-L}(C).}
\]
This remains true when the stored linear forms are arbitrary elements of $\F^b$, not necessarily codewords of $C$, and when continuation probes are adaptive.
\end{theorem}

\begin{proof}
Let $S$ be the row space of the persisted linear summary, with $\dim S\le L$.  Fix any reachable leaf of the continuation decision tree and let $J$ be the set of distinct source coordinates queried on its path.  All sources at this leaf have the same summary and the same values on $J$.  Exact output requires the affine leaf cell to be monochromatic under $W$, equivalently
\[
\ker U\cap\{x:x_J=0\}\subseteq\ker W.
\]
Taking orthogonal complements gives
\[
C\subseteq S+E_J,
\qquad E_J=\operatorname{span}\{e_j:j\in J\}.
\]
Reduce this inclusion modulo $E_J$.  The image of $C$ is contained in the image of $S$, so
\[
\dim(C/(C\cap E_J))\le\dim S\le L.
\]
Therefore $D=C\cap E_J$ has dimension at least $r-L$ and is supported inside $J$.  By the definition and monotonicity of generalized weights,
\[
|J|\ge d_{r-L}(C).
\]
This holds at every reachable leaf, so adaptivity cannot lower the worst-case depth.

Conversely, choose an $(r-L)$-dimensional subcode $D\le C$ attaining $d_{r-L}(C)$ and let $J=\supp D$.  Choose a linear complement $S\le C$ with $C=S\oplus D$ and $\dim S=L$.  Preprocessing stores a basis of the $S$-parities.  Continuation nonadaptively reads all coordinates in $J$, computes every parity in $D$, and reconstructs every row of $C$.  Hence $Q=d_{r-L}(C)$ is attainable.
\end{proof}

The proof is same-execution and permits arbitrary adaptive continuation logic.  It does not charge $Q$ twice for shared coordinates: all target rows reuse the same $J$ readings.

\subsection{Construction versus representation}

Given an attaining subcode $D$, a complement $S$ and the support $J$ can be verified and compiled by Gaussian elimination in polynomial time.  But obtaining an optimal $D$ is a separate construction problem.

\begin{corollary}[discovery firewall]\label{cor:hard}
For $L=r-1$,
\[
Q^*_{\rm lin}(r-1)=d_1(C),
\]
the minimum distance of $C$.  Thus computing the exact general linear-summary frontier, or always constructing an optimal endpoint, is NP-hard unless P$=$NP.
\end{corollary}

This corollary does not prove that executing a \emph{supplied} summary is hard.  It falsifies only the inference that a compact optimal representation entails polynomial discovery.  Structured families such as coordinate, repetition, Hamming, simplex, affine, or laminar codes may expose their attaining subcodes directly.

The complete ledger must therefore distinguish
\[
\mathfrak R=(P_{\rm discover},P_{\rm verify},L,Q,W_{\rm peak},R,H,D;\mathsf{access}).
\]
Theorem~\ref{th:frontier} determines only the optimum $L$--$Q$ slice after a suitable linear summary is chosen.  $P_{\rm discover}$ can dominate it.

\section{Contextual computational value}

For a current linear summary space $S$, define the independently specified continuation potential
\[
\Psi_C(S)=\min\{|J|:C\subseteq S+E_J\}.
\]
For a certified additional linear fact $u\cdot a$, write
\[
(u\cdot a)\mid_{v_\Psi},\qquad
v_\Psi(u\mid S)=\Psi_C(S)-\Psi_C(S+\langle u\rangle).
\]
The pipe denotes contextual value, not probability conditioning or divisibility.  This potential is monotone, so its raw value is nonnegative; a cost-aware net score may be negative after acquisition, verification, storage and application costs are charged.

The smallest XOR-share control already defeats additive intrinsic labels.  Let the target be $a_0$, and let $u=a_1$, $v=a_0+a_1$.  Then
\[
\Psi(0)=\Psi(\langle u\rangle)=\Psi(\langle v\rangle)=1,
\qquad
\Psi(\langle u,v\rangle)=0.
\]
Thus the two initial singleton values are $0,0$, while their joint value is $1$; after one share is known, the other has contextual value $1$.  Duplicates have value zero.

High value need not mean cheap acquisition.  For the length-$b$ repetition parity, $Psi(0)=b$ and storing the target parity gives value $b$, but direct construction from raw source bits itself reads all $b$ bits.  The full benefit is exactly offset by the naive acquisition work.  This is an easy affine control, not hardness.

The optimal envelope has marginal continuation savings
\[
Q^*_{\rm lin}(L-1)-Q^*_{\rm lin}(L)
=d_{r-L+1}(C)-d_{r-L}(C),
\]
which can exceed one.  These are contextual support savings, not entropy bits and not elapsed-time lower bounds.

\section{Nonlinear-summary red team}

The universal statement
\[
Q^*(L)=d_{r-L}(C)\quad\text{for arbitrary nonlinear summaries}
\]
is \status{OPEN}, not adopted.  A nonlinear summary fiber need not be an affine subspace, so the orthogonal-complement proof above does not apply.

The reproducer exhaustively tested every binary linear code of length $b\le4$, every one-bit summary map up to complement symmetry, and the exact optimal adaptive decision-tree depth on each summary fiber.  Across $60$ rank-at-least-two code cases and $1{,}672{,}200$ summary maps, no nonlinear advantage over $d_{r-1}(C)$ occurred.  For the $[4,3]$ even-weight code with hierarchy $(2,3,4)$, all $32{,}768$ one-bit summaries give optimum continuation depth at least $3$, matching the linear frontier.  This is \status{COMPUTATIONAL} evidence only.

An attempted maximal-domain enumeration at $b=5$ did not finish its first roughly $35$-second exploratory chunk and was terminated; it produced no evidence and is excluded from the claim.  Future work must use checkpointed SAT/set-cover search or a proof, rather than increasing a monolithic timeout.

\section{Peer referee findings}

\begin{longtable}{p{0.14\linewidth}p{0.61\linewidth}p{0.17\linewidth}}
\toprule
Artifact & Finding & Status\\
\midrule
ARC239 & The same-transcript flip proof for $p+Q\ge b$ is sound for all future parity masks.  Its fixed-read Fourier bound is compatible with this report but solves a different universal-query problem.  Its optimal-representation construction costs are explicitly polynomial only for the stated Hamming endpoints. & \status{PROVED/EXACT} restricted\\
SATURATION238 & The saturation argument correctly uses global Boolean degree-$Q$ extensions, Parseval equality and Walsh orthogonality.  The exact $(6,2,2)$ non-prefix point survives review. & \status{PROVED/EXACT} restricted\\
BATCH-LAMINAR175 & The finite-label tree-DP theorem provides a genuine structured polynomial discovery algorithm; it does not imply generalized-weight discovery for arbitrary codes or arbitrary SAT. & \status{PROVED/EXACT} restricted\\
BATCH-SWEEP237 & Sequential physical-head travel is a distinct ledger and cannot be substituted for indexed coordinate probes.  No newer predecessor report was verified this cycle. & \status{PROVED/EXACT} restricted\\
\bottomrule
\end{longtable}

ARC239 and the initial GitHub version of SATURATION238 lack the mandatory first report marker and final verification sentinel.  The later corrected Library copy is not treated as a GitHub readback.  The next GitHub publication should repair protocol markers without silently strengthening the mathematics.  BATCH-LAMINAR175 has both required markers.

\section{CORPUS REVIEWED}

Reviewed root \texttt{main.tex} (pre-run blob \texttt{775856b9c69defe2f36d3ea81a4e96521597a8b8}), \texttt{AGENTS.md}, \texttt{CODEX.md}, the full research skill, every current task specification, the newest verified MAIN content, the preceding LSPV-5 report and code/raw data, actual LSPV-1--4 sources materially bearing on parity summaries, the relevant governance constraints, primary coding/data-structure literature, and current GitHub implementations.  The root is substantive and has five balanced ownership regions; no placeholder text was used.

\section{DUPLICATION / CONTRADICTION AUDIT}

The rank-only lower bound and affine contextual potential in run211 predate this run.  The new exact object is the locality-sensitive hierarchy $d_{r-L}(C)$ under linear persistence, plus the construction-hardness firewall.  Wei's generalized weights themselves are a known coding invariant; no novelty claim is made for that concept.

There is no contradiction with ARC238/239 or SATURATION238: those require answering \emph{every} future parity mask, whereas Theorem~\ref{th:frontier} fixes one row code $C$.  There is no contradiction with LSPV-3's laminar planner: its menu structure is precisely extra information that can make discovery polynomial.

\section{WHAT CHANGED RELATIVE TO ROOT MAIN}

The root should refine the batch-rank statement in the linear-summary regime from a rank floor to the exact support frontier $Q=d_{r-L}(C)$, and immediately state that finding the attaining subcode is not generally free.  Rank is the number of independent answer bits; generalized support is the physical raw-coordinate locality.  Neither is a universal SAT runtime measure.

\section{ADOPTED / REJECTED / RETIRED PRIOR COMPONENTS}

\textbf{Adopted:} retained-source access; exact batch reuse; affine/XOR controls; source-dependent state accounting; independent discovery and execution ledgers; structured laminar planning.

\textbf{Rejected:} extrapolating the linear-subspace proof to arbitrary nonlinear summary fibers; using finite $b\le4$ evidence as an all-size theorem; treating an optimal description as cheaply findable.

\textbf{FALSIFIED/RETIRED:} ``shortest continuation representation implies polynomial summary construction'' for general linear codes, unless P$=$NP.  At $L=r-1$, exact optimal construction contains the NP-hard minimum-distance problem.

\section{STALE OR OVERSTATED CLAIMS FOUND}

Any claim that rank alone determines coordinate-query cost is stale: equal-rank codes can have different weight hierarchies.  Any use of $d_{r-L}$ for nonlinear summaries is currently overstated.  Any algorithmic upper bound that omits generalized-weight discovery is conditional on a supplied or structured attaining subcode.

\section{Reproducible computation}

The dependency-free Python checker finished in $20.4481$ seconds.  It verified $85$ binary code cases through length four, $9{,}439$ linear feasibility conditions, $60$ nonlinear one-bit-summary cases, $1{,}672{,}200$ summary maps, and $3{,}344{,}400$ cached restricted decision subproblems.  It found zero nonlinear advantages in that finite range and independently checked redundancy, XOR synergy, coordinate/even-weight/repetition hierarchies, and costly-high-value controls.  Code and raw JSON accompany this report.

\section{Ranked obligations and final claim}

\begin{enumerate}[leftmargin=*]
\item \textbf{OPEN, highest value:} prove the generalized-weight bound for arbitrary nonlinear fixed-length summaries, or produce a smallest explicit counterexample using checkpointed exact search.
\item \textbf{OPEN:} identify structured code/module classes for which an attaining subcode is uniformly discoverable in polynomial time and compose their $P,L,Q,W$ ledgers.
\item \textbf{OPEN:} authenticate a strict22 interface whose context batch realizes a nontrivial row code without granting an answer bit or affine solver for free.
\item \textbf{OPEN:} derive native sequential-tape and fixed-UTM counterparts charging head position, crossings, source rereads and symbol changes.
\end{enumerate}

\textbf{PROVED/EXACT, exact scope.}  For a fixed linear parity batch with row code $C$ of dimension $r$, a persisted $L$-bit linear summary followed by arbitrary adaptive indexed raw-coordinate probes has exact worst-case continuation optimum
\[
Q^*_{\rm lin}(L)=d_{r-L}(C).
\]
At $L=r-1$, constructing the general exact optimum contains minimum-distance computation; therefore compact continuation and polynomial construction are distinct obligations.  This is not a general SAT lower bound.

\textbf{Next specific lemma.}  Decide whether the same $d_{r-L}(C)$ lower bound survives arbitrary nonlinear fixed-length summaries.  Use a checkpointed exact cover/decision-domain formulation for the first unsettled length-five codes, or derive a fiber theorem strong enough to replace enumeration.

\begin{thebibliography}{9}
\bibitem{wei} V. K. Wei, ``Generalized Hamming Weights for Linear Codes,'' \emph{IEEE Transactions on Information Theory} 37(5):1412--1418, 1991. doi:10.1109/18.133259.
\bibitem{vardy} A. Vardy, ``The Intractability of Computing the Minimum Distance of a Code,'' \emph{IEEE Transactions on Information Theory} 43(6):1757--1766, 1997. doi:10.1109/18.641542.
\bibitem{ghwalg} R. San-Jos\'e et al., ``An Algorithm for Computing Generalized Hamming Weights and the Sage Package GHWs,'' \emph{ACM Transactions on Mathematical Software}, 2025/2026. doi:10.1145/3773284.
\bibitem{gal} A. G\'al and P. B. Miltersen, ``The Cell Probe Complexity of Succinct Data Structures,'' \emph{Theoretical Computer Science} 379(3):405--417, 2007. doi:10.1016/j.tcs.2007.02.047.
\end{thebibliography}

% END_OF_VERIFIED_REPORT::LSPV-5::GENERALIZED-SUPPORT-FRONTIER-REFEREE212
\end{document}
